%% mtest-example.tex -- example for the mtest class.
%%
%% Build with latexmk (two runs; see doc/NOTES.md).
\documentclass[onehalfspacing, 11pt, a4paper]{mtest}

\usepackage{mstuff}
\usepackage{msheet}

\title{Quadratische Funktionen}
\author{Marc Schlienger}
\date{\today}

%% \class appears in the head, \testnumber in the title table.
\class{Klasse 9b}
\testnumber{2}

%% Optional; both show up in the footer.
\license{\href{https://creativecommons.org/licenses/by/4.0}{CC BY 4.0}}

%%%%%%%%%%%%%%%%%%%%%%%%%%%%%%%%%%%%%%%%%%%%%%%%%%%%%%%%%%%%%%%%%%%%%%%%%%%%%%
\begin{document}

\maketitle

\mtestpart{Teil A -- ohne Hilfsmittel}

\msection[Scheitelpunktform]
Bringen Sie die folgenden Terme auf Scheitelpunktform. \mscore{6}
\begin{mexccolumns}(2)
	\exc $x^2 + 4x + 1$
	\exc $x^2 - 6x + 11$
\end{mexccolumns}

\msection[Nullstellen]
Bestimmen Sie die Nullstellen von $f(x) = x^2 - 5x + 6$. \mscore{4}

\msectionr[Zuordnung]
Ordnen Sie jedem Schaubild den passenden Term zu. \mscore{3}

\mtestpart{Teil B -- mit Hilfsmitteln}

\msection[Sachaufgabe]
Ein Ball wird senkrecht nach oben geworfen. Seine Höhe in Metern nach
$t$~Sekunden ist $h(t) = -5t^2 + 20t + 1$.
\begin{msectionlist}[start=4]
	\item Wann erreicht der Ball die maximale Höhe? \mscore{3}
	\item Wie hoch fliegt er? \mscore{2}
	\item Wann trifft er auf dem Boden auf? \mscore{3}
\end{msectionlist}

\begin{mbinfo}
	Runden Sie alle Ergebnisse auf zwei Nachkommastellen.
\end{mbinfo}

\mtotalscore

\mwish

\end{document}
